% ---------------------------------------------------------------------
%  CHƯƠNG 6 — KẾT LUẬN VÀ HƯỚNG PHÁT TRIỂN
% ---------------------------------------------------------------------
\chapter{Kết luận và hướng phát triển}
\label{ch:ketluan}

\begin{muctieu}
Chương này tổng kết phạm vi giải pháp, nêu giới hạn thực nghiệm, rút ra bài
học và đề xuất các hướng phát triển tiếp theo.
\end{muctieu}

\section{Kết luận}

Đồ án ``Xây dựng hệ thống gợi ý phối đồ'' xây dựng ứng dụng Android hỗ trợ
phân tích ảnh trang phục, gợi ý theo hồ sơ và bối cảnh, hỏi đáp, quản lý hồ sơ
cơ thể và thử trang phục theo tư thế camera. Giải pháp gồm Flutter, FastAPI,
YOLO, Gemini, MediaPipe Pose, Open-Meteo và Firebase. Firestore Rules đã dành
đường dẫn tủ đồ riêng cho từng UID, nhưng mã Flutter hiện chưa có luồng sử
dụng đường dẫn này và cũng chưa triển khai danh sách yêu thích; vì vậy hai
chức năng không được tính là kết quả đã hoàn thành.

Đối với thành phần nhận diện, đồ án đã huấn luyện checkpoint YOLO26n theo ba
vòng V1--V3 trên taxonomy 21 lớp. Phiên bản V3 kết hợp dữ liệu Roboflow và
DeepFashion, đạt kết quả validation tốt nhất mAP@0,5:0,95 là 36,69\% trong
lần chạy full; năm epoch fine-tune cuối đạt precision 55,18\%, recall
49,19\%, mAP@0,5 là 49,53\% và mAP@0,5:0,95 là 36,10\%. Các con số này xác
nhận quy trình huấn luyện hoạt động nhưng cũng cho thấy chất lượng chưa đủ để
kết luận đạt mục tiêu 80\% ban đầu.

Backend mặc định dùng \code{backend/models/best.pt}. Kiểm thử tự động ghi
nhận 104/104 test backend đạt; Flutter có 139 test đạt và 6 test widget lỗi.
Do chưa đo trên thiết bị Android thật, báo cáo chưa kết luận về mục tiêu thời
gian phản hồi không quá ba giây.

\section{Bài học rút ra}

Đề tài cho thấy chất lượng của một hệ thống trí tuệ nhân tạo không chỉ phụ
thuộc vào kiến trúc mô hình mà còn phụ thuộc dữ liệu, cách chia tập, quy
trình tiền xử lý và khả năng tích hợp với ứng dụng. Một mô hình có độ chính
xác cao nhưng phản hồi chậm hoặc khó triển khai vẫn có thể không phù hợp với
thiết bị di động.

Việc tách ứng dụng, API, dữ liệu và mô hình thành các thành phần rõ trách
nhiệm giúp kiểm thử và thay thế từng phần dễ hơn. Đồng thời, chức năng thử
đồ 2D minh hoạ giới hạn quan trọng của AI và xử lý ảnh: kết quả trực quan
không đồng nghĩa với mô phỏng chính xác kích cỡ hay đặc tính vật lý của
trang phục.

Quá trình V3 cũng cho thấy dữ liệu cần được kiểm tra trước huấn luyện: một số
nhãn cũ trộn định dạng detection và segmentation, vài hộp vượt miền chuẩn hoá
và phân bố lớp chênh lệch lớn. Việc chia DeepFashion theo \code{item\_id}
giúp tránh rò rỉ các góc chụp của cùng sản phẩm, nhưng ánh xạ giữa taxonomy
DeepFashion và 21 lớp đích vẫn tạo ra nhiễu nhãn.

\section{Hướng phát triển}

Một số hướng phát triển trực tiếp từ các hạn chế ở
Chương~\ref{ch:ketqua} gồm:
\begin{itemize}[leftmargin=1.6em, itemsep=3pt]
  \item mở rộng và cân bằng dữ liệu, bổ sung ảnh thực tế với nhiều nền, tư
        thế và điều kiện ánh sáng;
  \item nhận diện nhiều món đồ trong cùng ảnh và bổ sung thuộc tính như chất
        liệu, hoạ tiết, mùa hoặc hoàn cảnh sử dụng;
  \item cá nhân hoá gợi ý từ phản hồi, lịch sử lựa chọn và phong cách của
        từng người dùng;
  \item thay lớp ảnh 2D bằng phương pháp căn chỉnh theo tư thế hoặc mô hình
        thử đồ sinh ảnh, đồng thời đánh giá rủi ro sai lệch hình ảnh;
  \item hoàn thiện bảo mật, tối ưu thời gian phản hồi, bổ sung iOS và triển
        khai dịch vụ trên hạ tầng có khả năng mở rộng.
\end{itemize}

Mỗi hướng phát triển cần có dữ liệu, tiêu chí đánh giá và kế hoạch kiểm thử
riêng; không nên coi việc thay mô hình mạnh hơn là đủ để giải quyết mọi hạn
chế của hệ thống.
